\begin{statement}{P3K.31}{Proved}
For every actual family and positive parameter loop of P3K.27 or
P3K.30, retain its physical delta cocharacter and all exterior
gluings, and use any allowed index-one periodic unimodular cusp
fan \(T(v,w)\). Write
\[
 \delta_{\rm circle}=r v+s w,\qquad
 S=r^2-rs+s^2,\qquad \gcd(r,s)=1.
\]
Let \(k,\ell,A=(4k-\ell)/24\) be the initial vertical-edge
integers in the actual Cox framing of those statements.
The oriented return class for the specified fan is
\[
 [\Sigma_{\Phi,T}]
    =A+\frac{\ell}{8}(S-1)-\frac{k}{8}(S^2-1)
        \pmod{28}.                                      \tag{P3K.31.0}
\]
The expression is integral. It is computed by the actual common
exterior, stable tangent-frame change, and complex normal reset
through every periodic flop.

For K27 this becomes
\[
           4-3S+\frac d8(3S^2-7)\pmod{28},
       \qquad d=8,14,20,30.
\]
For K30, with its proved \(H=Ma^2/d\) and page Euler characteristic
\(\chi\), it becomes
\[
                 \chi(4-3S)+\frac H8(3S^2-7)\pmod{28}.
\]
The exact finite ranges for K27 are given in the proof.
 The original construction's single \(w\)-loop is treated separately
 in P3K.33 by its explicit simultaneous endpoint descent. No cover
 class is divided.
\end{statement}
\begin{proof}
\par\medskip\noindent\textbf{1. The three exact objects that must be transported.}\enspace

Keep K29's positive complex normal meridians \(m_i\) and relative
dual cycles \(m_j\cdot D_i=\delta_{ji}\). Let \(V\) denote the
variation matrix of the actual relative return
\(\operatorname{Var}_F(D)=F(D)-D\). Choose a global stable tangent
frame of the actual six-sphere and restrict it to the common exterior.
The derivative of \(F\), equal to the identity near its boundary,
lifts to the stable Spin group. Write
\[
 \delta_i=\langle c(F)^*\eta,D_i\rangle,\qquad
       \delta=(\delta_0,\delta_1,\delta_2)^T,
                                                        \tag{P3K.31.1}
\]
where the integral third generator \(\eta\) is normalized so that
the stable map \(q\mapsto L_q\), left quaternion multiplication,
has degree \(+1\).

The third object is the oriented central return class \(s\in\mathbb Z/28\).
The initial delta-retaining family has K23 relative integers
\((m,n)\), with
\[
       V=\operatorname{diag}(v,0,0),\quad
       v=-(m+n)=k_{\rm Jin},\quad
       \delta=(n-m,0,0)^T,\quad s=A.
                                                        \tag{P3K.31.2}
\]
The formula for \(\delta\) will follow from the literal product-frame
comparison below; it is not inferred from an absolute Pontryagin class.

A different global stable frame gives a map from the whole six-sphere
to Spin. Its third class is zero because \(H^3(S^6)=0\).
K29's exact frame-change pairing therefore shows that (P3K.31.1) is
independent of that global choice at the degree being used.
No assertion about a higher framing homotopy group is needed.

\par\medskip\noindent\textbf{2. The first-resolution product differs by one normal generator.}\enspace

Use the actual first small resolution and the K23 quaternion product
\[
 u=a_0+j a_1,\qquad
 c=c_0+j\overline{c_1},\qquad
 q=uc,\qquad h=u i u^{-1}.
                                                        \tag{P3K.31.3}
\]
The positive complex meridian is the \(q\)-sphere. The base
\(h\)-sphere has the opposite of the complex curve orientation.
The exterior product is oriented \(D^3_h\times S^3_q\).
At a fixed base point, an outward radial vector of \(D^3_h\)
corresponds to the negative of the outward normal radial
vector \(q\) of the curve tube: the two are the collar
directions on opposite sides of their common boundary.

The derivative has a source-to-normal shear that is
retained before making the following frame comparison.
The inverse product map is literally \(c=u^{-1}q\).
At fixed \(q\), a source derivative therefore gives
\[
                   dc=-u^{-1}(du)u^{-1}q.
\]
Together with the source derivative and the radial
normal column, the full frame matrix has block form
\[
       \begin{pmatrix}
         B_2&0\\ A(q)&L_qR_4
       \end{pmatrix}.
\]
Here \(A(q)\) is precisely the normal component just
displayed, after the fixed source basis change.
Every matrix
\[
       \begin{pmatrix}
         B_2&0\\ tA(q)&L_qR_4
       \end{pmatrix},\qquad 0\le t\le1,
\]
is invertible and has the same positive determinant.
This is the actual continuous triangular homotopy
removing that derivative term. In the second product
the inverse is \(u=qc^{-1}\), so at fixed \(q\)
\(du=-qc^{-1}(dc)c^{-1}\); the identical block homotopy,
with the fixed reflection \(J\) in its normal block,
retains and removes its corresponding shear.

A fixed linear change of the three Cartesian \(h\)-vectors
therefore gives the two tangent vectors of the curve and
the vector \(-q\). The tangent curve pair has orientation
opposite to its complex orientation. The three left
quaternion vectors of the normal sphere are \(qi,qj,qk\).
Thus, in a local complex tangent frame of the ambient
six-manifold, the actual product frame along the meridian
has the form
\[
       \operatorname{diag}(B_2,L_q R_4),\qquad
       \det B_2=-1,\quad
       R_4=\operatorname{diag}(-1,1,1,1).
                                                        \tag{P3K.31.4}
\]
Both determinants are negative, so the total matrix is
orientation preserving. After multiplying by its constant
value at \(q=1\), its class is exactly
\(\operatorname{diag}(I_2,L_q)\).
The ambient stable frame extends over the small normal
four-disc, so its restriction has zero third class.
Consequently the product frame differs from the ambient
frame by degree \(+1\) on its positive complex meridian.
It has degree zero on the other two meridians, since its
frame extends over their normal discs in the selected
single-curve exterior.

Let \(g_i\) be this change from ambient to the selected
first-resolution product frame. Its absolute third
cohomology vector in the common exterior is exactly
\[
                              [g_i]=e_i.
                                                        \tag{P3K.31.5}
\]
By K29, the actual derivative in the product frame
differs from the ambient derivative by
\[
           -\langle[g_i],\operatorname{Var}_F(D_i)\rangle=-V_{ii}.
\]
The ordered product disc and the relative dual
disc have opposite signs:
\[
              D_i=-[D^3_h],\qquad
              m_i\cdot[D^3_h]=-1.
                                                        \tag{P3K.31.6}
\]
K23's derivative on \([D^3_h]\) is \(-2n_i\). Equations
(P3K.31.5)--(P3K.31.6) give the exact first-resolution formula
\[
                    2n_i=\delta_i-V_{ii},\qquad
                    m_i^{\rm rel}=-V_{ii}-n_i.
                                                        \tag{P3K.31.7}
\]
Here \(m_i^{\rm rel}\) is a relative group integer, not a
normal meridian. In the initial fan K29 makes all three
quantities zero for \(i=1,2\), and the fixed edge has
\(\delta_0=v+2n=n-m\), proving (P3K.31.2).

\par\medskip\noindent\textbf{3. The literal ambient-frame change under the flop is \(+2e_i\).}\enspace

First keep the unchanged \(q\) coordinate of the actual
two small resolutions. The second product has coordinates
\[
             k=c^{-1}ic,\qquad q=uc,
\]
and the common boundary map is
\[
                 T(h,q)=(q^{-1}h q,q).
                                                        \tag{P3K.31.8}
\]
K23 proves that this is the literal actual-family map,
independent of the parameter angle, including odd Cox
loops and the transported bounded-unit truncation collars.
Its displayed extension over \(D^3\times S^3\) gives a
normal collar representative of the same boundary map.

The derivative in Cartesian base and left quaternion
normal frames is triangular. Its base block is
\(\operatorname{Ad}_{q^{-1}}\), its normal tangent block
is the identity, and its off-diagonal block contracts
linearly to zero. Thus its stable third class on the
old positive \(q\)-meridian is
\[
                         [c_T]=-2.
                                                        \tag{P3K.31.9}
\]
This sign and factor follow from the actual quaternion
action: \(\operatorname{Ad}_q=L_qR_{q^{-1}}\) has stable
degree \(+2\) in the normalization (P3K.31.1), while inversion
of \(q\) reverses its degree. Equivalently the marked
adjoint clutch has \(p_1=-4c_2\), with the same direct
orientation calculation as K23; the left quaternion
clutch is its degree-one stable normal generator.

In the second product, the \(k\)-sphere preserves the
complex curve orientation and the unchanged \(q\)
coordinate reverses the complex normal orientation.
Let \(J\) be the fixed reflection of that normal real
four-plane. The analogue of (P3K.31.4), parametrized by
the unchanged \(q\)-sphere, is
\(\operatorname{diag}(B_2',J L_qR_4)\), with
\(\det B_2'=+1\).
After its constant value is removed its stable class
is \(J L_qJ^{-1}\), still degree \(+1\).
Indeed conjugation by \(J\) becomes conjugation by
\(\operatorname{diag}(J,-1)\) after one stabilization;
the latter lies in the connected special orthogonal
group and hence induces the identity on the third
homotopy group. The new positive complex meridian is
\(-[S^3_q]\), so this same product-frame difference has
degree \(-1\) on the new complex meridian, exactly as
required by the old/new normal orientations.

Write \(P_o,P_n\) for the two product frames and
\(\sigma_o,\sigma_n\) for the two global ambient stable
frames. Put \(P_o=\sigma_og_o\), \(P_n=\sigma_ng_n\).
Pulling the new frame back by the actual map gives
\[
 (dT)^{-1}\sigma_n
       =\sigma_o\,g_o\,c_T^{-1}\,(g_n\circ T)^{-1}.
                                                        \tag{P3K.31.10}
\]
On the old positive complex meridian the three displayed
degrees are \(+1\), \(-(-2)\), and \(-1\), respectively.
The exact ambient-frame change is therefore
\[
                         [G]=1+2-1=2.
                                                        \tag{P3K.31.11}
\]
On either other curve's normal meridian the birational
map is a smooth complex isomorphism on the whole normal
four-disc, so the change has degree zero. Hence for a
flop of edge \(i\), its absolute common-exterior third
class is exactly \(2e_i\). This computes the relative
normal Thom correction; it is not chosen from the
vanishing of an absolute characteristic class.

\par\medskip\noindent\textbf{4. Resetting the first-resolution complex convention.}\enspace

The unchanged-\(q\) second product in (P3K.31.8) is the one
used in K23's conjugation formula
\[
             (m,n,s)\longmapsto(-n,m+2n,s+n).
                                                        \tag{P3K.31.12}
\]
For the next selected edge we must return to the first
product whose normal quaternion preserves its complex
orientation. This reset has an explicit map.

Write a quaternion as \(x_0+ix_1+jx_2+kx_3\), and set
\[
 J(x_0+ix_1+jx_2+kx_3)=x_0+ix_1+jx_2-kx_3.
                                                        \tag{P3K.31.13}
\]
It is the anti-automorphism \(J(q)=k\bar q k^{-1}\).
For the new first resolution its base quaternion is
\(J(c)=c_0+jc_1\), its normal quaternion is
\(J(u)=a_0+j\bar a_1\), and their product is
\[
                      q_{\rm new}=J(c)J(u)=J(q).
                                                        \tag{P3K.31.14}
\]
Its Hopf base is \(J(k)\). Thus the exact reset is
\[
                       U(y,q)=(J(y),J(q)).
                                                        \tag{P3K.31.15}
\]
The base three-reflection has determinant \(-1\) on
\(\operatorname{Im}\mathbb H\), and the normal
four-reflection has determinant \(-1\) on
\(\mathbb H\). Their product is \(+1\), so the map
preserves both the six-dimensional exterior and
the seven-dimensional filling orientation.
The map extends over \(D^3\times D^4\).
For K23's conjugation-equivariant radial \(\zeta\),
\(J\zeta(Jy)=\zeta(y)\), so literal substitution gives
\[
                    ULU^{-1}=R,\qquad URU^{-1}=L.
                                                        \tag{P3K.31.16}
\]
Its orientation-preserving extension fixes the central
disc-supported class. Hence the selected first-convention
pair after the flop is
\[
                        (m+2n,-n),\qquad s'=s+n.
                                                        \tag{P3K.31.17}
\]
In particular the reverse flop changes the central class
by \(-n\), and restores the original class.
This reset also retains the actual parameter family.
For the Cox product action \(q\mapsto a_\theta q b_\theta\),
the same literal anti-automorphism gives
\(q_{\rm new}\mapsto J(b_\theta)q_{\rm new}J(a_\theta)\).
The corresponding base action follows by substituting
in its Hopf coordinate. The map \(U\) is independent of
\(\theta\); it therefore intertwines the two existing
family products, including their endpoint signs, and
does not add a parameter-dependent frame loop.

\par\medskip\noindent\textbf{5. The complete common-exterior recurrence.}\enspace

Let \(M_i\) be K29's actual meridian matrix: its \(i\)-th
column is \(-e_i+2\sum_{j\ne i}e_j\), and its other
columns are unchanged. The relative basis change is
\(Q_i=M_i^{-T}\). Equations (P3K.31.10)--(P3K.31.11), the exact
frame-change pairing, and the same actual common-exterior
return give
\[
 \begin{split}
 V'&=M_iVM_i^T,\\
 \delta'&=M_i(\delta-2Ve_i),\\
 s'&=s+\frac{\delta_i-V_{ii}}2\pmod{28}.
 \end{split}                                             \tag{P3K.31.18}
\]
All quantities use the first-resolution complex normal
convention fixed explicitly in Section 2.
The first line is K29. For the second, changing ambient
frame subtracts \(V^T(2e_i)=2Ve_i\) on the old relative
basis, and evaluating on \(Q_i\) multiplies by
\(Q_i^T=M_i\). The final line is K23 with (P3K.31.7).
The literal reset (P3K.31.15) shows that this last sign is
the one for a further actual flop, not an unchanged-\(q\)
second-product integer.

The initial state (P3K.31.2) gives after a fixed-edge flop
\[
 \begin{split}
 V'&=v\begin{pmatrix}1&-2&-2\\-2&4&4\\-2&4&4\end{pmatrix},\\
 \delta'&=(m+3n)(-1,2,2)^T,\\
 n'_0&=-n,\qquad n'_1=n'_2=3m+5n .
 \end{split}                                             \tag{P3K.31.19}
\]
The selected reverse-flop value and the two other
values are different because their actual relative
cycles are the different columns in K29. The latter
retain both original free relative integers.

\par\medskip\noindent\textbf{6. Closed formula in the actual tangent weights.}\enspace

We now derive the closed formula from the actual recurrence.
Write the initial coefficient as \(v=-(m+n)\). Along a
sequence of actual periodic flips, write
\[
                  V=vww^T,\qquad \delta=a w,
 \qquad w_{\rm initial}=(1,0,0),\quad a_{\rm initial}=n-m .
\]
Equation (P3K.31.18) gives
\[
       w'=M_iw,\qquad a'=a-2v w_i.
                                                        \tag{P3K.31.20}
\]
For \(S=w_0+w_1+w_2\), the same matrix gives
\(S'=S+2w_i\). Consequently \(a+vS=-2m\) is constant.
The exact central increment is
\[
 n_i=\frac{-2m\,w_i-vw_i(S+w_i)}2
     =-\frac m2(S'-S)-\frac v8(S'^2-S^2).
                                                        \tag{P3K.31.21}
\]
It telescopes to
\[
          s(S)=s(1)-\frac m2(S-1)
                         +\frac{m+n}{8}(S^2-1)\pmod{28}.
                                                        \tag{P3K.31.22}
\]

These integers have a literal local weight interpretation.
At a chosen triple point let the physical delta tangent
weights along the three complex double curves be
\((\lambda_0,\lambda_1,\lambda_2)\), with sum zero.
The smooth toric chart and primitive delta lattice vector
give \(\gcd(\lambda_0,\lambda_1,\lambda_2)=1\).
Put
\[
              w_i=-\lambda_j\lambda_k\quad
                     (\{i,j,k\}=\{0,1,2\}).
                                                        \tag{P3K.31.23}
\]
The initial fixed-edge weights \((0,1,-1)\) give
\((1,0,0)\). In the first small-resolution chart the
coordinates have weights \((\lambda_i,\lambda_j,\lambda_k)\).
On the second resolution's \(c_j\ne0\) chart, the actual
coordinates are \(c_k/c_j,a_0c_j,a_1c_j\), of weights
\[
                  (\lambda_k-\lambda_j,\lambda_j,-\lambda_k).
                                                        \tag{P3K.31.24}
\]
Substitution in (P3K.31.23) gives exactly \(w'=M_iw\).
At the other triple point all weights are negated and
permuted, so the sum below is unchanged. Thus
\[
             S=-\sum_{j<k}\lambda_j\lambda_k
               =\frac{\lambda_0^2+\lambda_1^2+\lambda_2^2}{2}
               =p^2+pq+q^2
                                                        \tag{P3K.31.25}
\]
for the actual primitive weight pair
\((\lambda_0,\lambda_1,\lambda_2)=(p+q,-p,-q)\).
In particular \(S\) is positive and odd, so the fractions
in (P3K.31.22) are integral.

For the primitive K22--27 families, \(m=6\), \(m+n=3d\),
and \(s(1)=1-d/2\). The resulting expression is
\[
        s(S)=1-\frac d2-3(S-1)+\frac{3d}{8}(S^2-1)
             =4-3S+\frac d8(3S^2-7)\pmod{28}.
                                                        \tag{P3K.31.26}
\]
For \(S=1,3,7\) it gives, respectively, the initial
delta-retaining fan, its fixed-edge flop, and a further
flop of either other edge. The second case is exactly
K23's \(-5A\). The actual frame and periodic flop morphisms above prove these values.

\par\medskip\noindent\textbf{7. Actual fan coordinates and the finite residue computation.}\enspace

P3F.35 supplies the exact global scope of the iteration.
Every allowed one-vertex unimodular index-one fan is
\(T(v,w)\) for an integral basis, and its explicit
periodic diagonal moves connect it to the vertical-edge
fan. The common blowups and contractions commute with
the actual period lattice, preserve the compact
six-sphere, and have no nontrivial stabilizer of an
individual selected edge cell. These are exactly the
local and global maps used in (P3K.31.8)--(P3K.31.18).

Write the physical cocharacter as
\(\delta=a v+b w\). In the triangle with rays
\((0,1),(v,1),(v+w,1)\), its height-zero coefficients
are \((-a,a-b,b)\). They are the three actual
coordinate weights, because the rays are an integral
basis of the three-dimensional toric lattice.
Consequently the integer in (P3K.31.25), expressed in the
marked fan rather than a chosen flip word, is
\[
                         S=a^2-ab+b^2.
                                                        \tag{P3K.31.27}
\]
The pair \((a,b)\) is primitive. Conversely every
primitive pair occurs for an actual integral basis:
choose a second lattice vector \(\eta\) with
\(\det(\delta,\eta)=1\), and integers \(c,d\) with
\(ad-bc=1\). The literal matrix
\[
       (v\ \ w)=(\delta\ \ \eta)
                        \begin{pmatrix}a&c\\b&d\end{pmatrix}^{-1}
\]
has determinant one and sends \((a,b)^T\) to \(\delta\).
It gives the fan \(T(v,w)\) of P3F.35 in the actual
period lattice.
Thus (P3K.31.27), together with (P3K.31.22), assigns the
derived value to each actual marked fan
without identifying distinct fans merely from an
abstract triangulation.

The arithmetic range is a finite exact calculation.
For arbitrary integers \(m,t=m+n\) and odd \(S\),
the expression
\[
                  -m(S-1)/2+t(S^2-1)/8
\]
is integral. Replacing \(S\) by \(S+112\) changes it
by \(28(-2m+tS+56t)\). Hence its value modulo28 is
determined by \((a,b)\) modulo112.
A residue pair admits a primitive integral lift
exactly when \(\gcd(a,b,112)=1\). Necessity is
immediate. For sufficiency choose a nonzero lift
of \(b\). A prime dividing this lift and 112 does
not divide the given residue of \(a\). For each
of its other, finitely many prime divisors, choose
the multiple of112 added to \(a\) to avoid zero
modulo that prime. The ordinary finite Chinese
remainder map supplies those choices simultaneously.
The resulting lifts are coprime.

In fact every odd residue of \(S\) modulo112 occurs.
Modulo16 take \(p=1\) and \(q=0,\ldots,7\); the
values \(p^2+pq+q^2\) are
\(1,3,7,13,5,15,11,9\), all odd residues.
Modulo7 the nonzero pairs
\[
 (1,2),(1,0),(0,3),(1,1),(2,0),(2,2),(1,3)
\]
give residues \(0,1,2,3,4,5,6\), respectively.
The ordinary Chinese remainder map combines these
choices, with a primitive residue pair modulo112;
the preceding explicit lift then makes it primitive
over the integers. Thus the finite calculation can
equally be checked on all56 odd residues of \(S\).

The executed certificate enumerates precisely those
primitive residue pairs, with witnesses for every
resulting class. For the four K22 denominators it gives:

\begin{center}\begin{tabular}{c p{0.78\linewidth}}\toprule
\(d\) & Derived classes over all primitive weight residues\\
\midrule
\(8\) & \(1,3,5,11,15,17,19,25\)\\
\(14\) & \(0,2,4,6,8,10,12,14,16,18,20,22,24,26\)\\
\(20\) & \(3,5,7,11,17,19,21,25\)\\
\(30\) & \(0,4,6,8,14,18,20,22\)\\
\bottomrule\end{tabular}\end{center}

These are the exact arithmetic ranges of the proved
geometric formula. Their executable
source, complete JSON and captured output are
\path{code/part3_fourfold/common_exterior_frame_recurrence.py},
\path{code/part3_fourfold/common_exterior_frame_recurrence.json},
and \path{code/part3_fourfold/common_exterior_frame_recurrence.out}.
The certificate checks the literal anti-quaternion
identity, the orientation reset, all three recurrence
involutions, the actual toric weight maps, and the
telescoping potential before evaluating the residues.

\par\medskip\noindent\textbf{8. The separately proved original-parameter descent.}\enspace

The initial simultaneous three-curve product in the present
statement is the actual K27 scalar-loop product or K30
one-gamma product. P3K.33 separately constructs the original
single positive \(w\)-loop and every higher \(w_h\)-loop
by an integral endpoint cochain. Its bounded origin closes
all three cusp charts and both triple points simultaneously;
the explicit physical-delta adjustment fixes the Hopf normal
vector. Thus P3K.33 supplies the exact single-loop product
needed to apply the recurrence proved here, for both sphere
determinant signs. It obtains
\[
             KL\frac{(S-1)(S-3)}8\pmod{28}
\]
directly from that product and its individual integer pair.
No class in \(\Theta_7\) is divided by the cover degree.
\end{proof}
